\chapter{Counting principles}\label{ch:count}
\index{counting}

Combinatorics is arithmetic with a story attached. Every formula in this
chapter --- $\binom{n}{k}$, $n!$, the stars-and-bars count --- is a count of
\emph{something concrete}, and the fastest proofs are the ones that name
the object and count it two ways. After this chapter you will be able to
read a counting problem, name the underlying set, pick the right model
(ordered vs unordered, distinct vs repeated), and verify small cases
against brute force before trusting a closed form.

% from-zero:begin
\begin{rulebox}[Before you start]
Chapters~2--4. No number theory is needed here. \emph{Recap:} counting means finding the size of a finite set without listing it; $n!=1\cdot2\cdots n$ is the number of ways to order $n$ different things.
\end{rulebox}
\begin{wins}[After this chapter you can]
\begin{itemize}
\item choose between the sum and product rules and know when each applies;
\item count arrangements, selections and distributions with factorials, binomials and stars-and-bars;
\item prove a counting identity by a bijection or a committee argument;
\item use the pigeonhole principle to prove that something must exist.
\end{itemize}
\end{wins}
\begin{principle}[How to think about this]
Counting is about structure, not formulas. Before reaching for $\binom{n}{k}$, ask: ``what are the choices, in what order do I make them, and do different choice-paths ever produce the same object?'' If you can describe how to \emph{build} each object exactly once, you have the count.
\end{principle}
% from-zero:end

\section{Sum and product rules}\label{sec:rules}
\index{sum rule}\index{product rule}

Two rules generate every formula in the chapter:
\begin{itemize}
\item \textbf{Sum rule.} If the objects split into disjoint cases of sizes
$a$ and $b$, there are $a + b$ objects. ``Disjoint'' is the licence ---
shared cases get double-counted.
\item \textbf{Product rule.} If an object is built by a sequence of
independent choices of sizes $a_1, a_2, \dots, a_k$, there are
$a_1 a_2 \cdots a_k$ objects --- one per root-to-leaf path in
Figure~\ref{fig:dectree}'s decision tree.
\end{itemize}

\begin{figure}[ht]\centering
% Figure: decision tree for the product rule --- 3 x 2 x 2 = 12 outcomes.
% Source for Figure~\ref{fig:dectree}. TikZ only; regenerable by compiling main.tex.
\begin{tikzpicture}[
  nd/.style={circle,draw=rulegray,minimum size=3.2mm,inner sep=0pt},
  lf/.style={rectangle,draw=accent,inner sep=1.6pt,minimum size=3mm},
  lab/.style={font=\footnotesize\sffamily,text=slate},
  ed/.style={draw=rulegray},
]
\node[nd] (r) at (0,0) {};
% level 1 nodes at y = +2,0,-2 ; level 2 at +-0.7 offsets ; leaves +-0.28
\foreach \i/\y in {0/2.0,1/0.0,2/-2.0}
{
  \node[nd] (a\i) at (1.7,\y) {};
  \draw[ed] (r) -- (a\i);
  \foreach \j/\dy in {0/0.7,1/-0.7}
  {
    \pgfmathsetmacro{\by}{\y+\dy}
    \node[nd] (b\i\j) at (3.4,\by) {};
    \draw[ed] (a\i) -- (b\i\j);
    \foreach \k/\ddy in {0/0.28,1/-0.28}
    {
      \pgfmathsetmacro{\ly}{\by+\ddy}
      \node[lf] at (5.1,\ly) {};
      \draw[ed] (b\i\j) -- (5.1,\ly);
    }
  }
}
\node[lab] at (0,2.6) {choice 1: 3 ways};
\node[lab] at (1.9,2.6) {choice 2: 2 ways};
\node[lab] at (4.4,2.6) {choice 3: 2 ways};
\node[lab] at (2.6,-2.9) {$3 \times 2 \times 2 = 12$ leaves --- every root-to-leaf path is one outcome};
\end{tikzpicture}

\caption{The product rule as a decision tree. Notice the multiplication
is in the \emph{shape}: each level doubles the fan-out, and leaves are
outcomes --- not choices. Counting \emph{paths}, not \emph{nodes}, is the
discipline.}
\label{fig:dectree}
\end{figure}

The mistakes are symmetric: multiplying when cases should be added
(``3 shirts or 4 trousers'' is $3+4$, not $3\cdot4$), and adding where a
build-process should multiply (``3 shirts then 4 trousers'' is an outfit
count, $3\cdot4$). The diagnostic is always: is the second choice
\emph{conditional on} the first?

\section{Permutations and factorials}\label{sec:perms}
\index{permutation}\index{factorial}

Orderings of $n$ distinct objects: $n!$ --- the product rule applied
$n$ times ($n$ choices for first, $n-1$ for second, \dots). A
$k$-permutation (ordered selection of $k$ from $n$) is
$n!/(n-k)!$: the last $n-k$ slots ``don't matter'', and dividing by
their arrangements is honest cancellation, not a trick.

\begin{example}\label{ex:mississippi}
How many distinct arrangements does \textsc{mississippi} have?
\par\smallskip
\textbf{Strategy.} Count as if all $11$ letters were distinct, then undo
the overcount: identical letters can't be told apart, so each arrangement
was counted once per internal shuffle of every repeated letter.
\par\smallskip
\begin{solution}
$11$ letters: one M, four I's, four S's, two P's. Naive count $11!$
treats the I's as distinguishable $I_1, \dots, I_4$ --- each real
arrangement appears $4!$ times (I-shuffles) $\times$ $4!$ (S-shuffles)
$\times$ $2!$ (P-shuffles). So the count is
$\dfrac{11!}{4!\,4!\,2!} = 34\,650$. The same division-by-shuffles gives
the multinomial formula $\binom{n}{n_1,\dots,n_k}$ in general.
\end{solution}
\end{example}

\begin{exercise}[easy]\label{exr:perm}
How many ways to award gold, silver, bronze in a 20-runner race? What if
ties were allowed --- does the formula still apply?
\end{exercise}

\section{Binomial coefficients and Pascal's triangle}\label{sec:binom}
\index{binomial coefficient}\index{Pascal's triangle}

\begin{definition}
$\binom{n}{k}$ (``$n$ choose $k$'') is the number of $k$-element subsets
of an $n$-element set: $\binom{n}{k} = \frac{n!}{k!\,(n-k)!}$.
\end{definition}

Pascal's triangle builds the table by the rule
$\binom{n}{k} = \binom{n-1}{k-1} + \binom{n-1}{k}$ --- itself a sum-rule
proof (the subsets either contain element $n$ or don't).
Figure~\ref{fig:pascal} highlights the hockey-stick identity, a workhorse
for summing columns.

\begin{figure}[ht]\centering
% Figure: Pascal's triangle, hockey-stick path highlighted.
% Source for Figure~\ref{fig:pascal}. TikZ only; regenerable by compiling main.tex.
% Highlighted cells (stick): (r,c) = (2,2),(3,2),(4,2),(5,2) + answer (6,3).
\begin{tikzpicture}[
  cell/.style={font=\small},
  hl/.style={font=\small\bfseries,text=accent},
  lab/.style={font=\footnotesize\sffamily,text=slate},
]
\foreach \r/\vals in {0/{1},1/{1,1},2/{1,2,1},3/{1,3,3,1},4/{1,4,6,4,1},5/{1,5,10,10,5,1},6/{1,6,15,20,15,6,1}}
{
  \foreach [count=\c from 0] \v in \vals
  {
    \pgfmathsetmacro{\px}{\c*0.7 - \r*0.35}
    \node[cell] at (\px,-\r*0.55) {\v};
  }
}
% redraw highlighted cells in accent (stick c=2 rows 2-5, answer r6c3)
\foreach \r/\c/\v in {2/2/1,3/2/3,4/2/6,5/2/10,6/3/20}
{
  \pgfmathsetmacro{\px}{\c*0.7 - \r*0.35}
  \node[hl] at (\px,-\r*0.55) {\v};
}
\node[lab,anchor=west] at (2.1,-1.1) {stick: $1+3+6+10$};
\node[lab,anchor=west] at (2.1,-1.5) {$= 20 = \binom{6}{3}$};
\node[lab] at (0,-4.0) {$\binom{2}{2}+\binom{3}{2}+\binom{4}{2}+\binom{5}{2}=\binom{6}{3}$:
 a column of terms sums into the cell one step right-down};
\end{tikzpicture}

\caption{Hockey-stick in Pascal's triangle. Notice the highlighted stick
$1+3+6+10$ summing to the cell one step right-down, $20$: the identity
$\sum_{i=r}^{n}\binom{i}{r} = \binom{n+1}{r+1}$ is the triangle's own
shape.}
\label{fig:pascal}
\end{figure}

\section{The binomial theorem, proved two ways}\label{sec:binthm}
\index{binomial theorem}

\begin{theorem}[Binomial theorem]
$(x + y)^n = \sum_{k=0}^{n} \binom{n}{k} x^k y^{n-k}$.
\end{theorem}

\emph{Proof 1 (committee).} Expanding $(x+y)^n$ is a product-rule walk:
from each of the $n$ factors pick $x$ or $y$. The term $x^k y^{n-k}$
appears once for every way to choose which $k$ factors gave $x$ --- that
is the definition of $\binom{n}{k}$.
\emph{Proof 2 (induction).} Multiply $(x+y)^{n-1}$ by $(x+y)$ and collect
$x^k y^{n-k}$: it arrives from $x \cdot x^{k-1}y^{n-k}$ (coefficient
$\binom{n-1}{k-1}$) plus $y \cdot x^k y^{n-1-k}$ ($\binom{n-1}{k}$), and
Pascal's rule completes it.

Setting $x = y = 1$ gives $\sum_k \binom{n}{k} = 2^n$ --- the next
example's committee argument proves the same identity directly.

\begin{example}\label{ex:committee}
Prove $\sum_{k=0}^{n} \binom{n}{k} = 2^n$ by counting in two ways.
\par\smallskip
\textbf{Strategy.} Count the same object by two different routes and set
the counts equal: a subset of $[n]$ can be counted \emph{by size}
($\binom{n}{k}$ per size, summed) or \emph{by membership} (each element
independently in or out: $2^n$). Same set, two formulas.
\par\smallskip
\begin{solution}
The left side counts subsets of $[n]$ partitioned by cardinality:
$\binom{n}{0} + \binom{n}{1} + \cdots + \binom{n}{n}$. The right side
counts the same subsets by the product rule: $n$ independent
in/out decisions, $2^n$ outcomes. Equal by definition of ``counting the
same thing'' --- that \emph{is} the proof.
\end{solution}
\end{example}

\begin{exercise}[medium]\label{exr:committee}
Prove $\sum_{k} k\binom{n}{k} = n\,2^{n-1}$ the same way (count
``subsets with a marked element'' two ways).
\end{exercise}

\section{Bijection is counting}\label{sec:biject}
\index{bijection!counting by}

Chapter~4's bijections pay rent here: a bijection between two finite sets
makes their sizes equal, so \emph{to count $A$, count an easier $B$}.
That is why ``counting in two ways'' proves identities and why renaming
the object is the first move on a hard problem.

\begin{example}\label{ex:lattice}
Count lattice paths on an $m \times n$ grid: paths from $(0,0)$ to
$(m,n)$ using only right-steps R and up-steps U.
\par\smallskip
\textbf{Strategy.} Name the bijection: a path is a word over
$\{R, U\}$ with exactly $m$ R's and $n$ U's --- length $m+n$, and the
path is determined by which positions are R.
\par\smallskip
\begin{solution}
Paths $\leftrightarrow$ words of length $m+n$ with $m$ R's is a
bijection (Figure~\ref{fig:lattice}). Choosing which $m$ of the $m+n$
positions are R determines the word, so the count is $\binom{m+n}{m}$.
\end{solution}
\end{example}

\begin{figure}[ht]\centering
% Figure: lattice-path bijection --- a 4x3 grid path vs the word RRURURR.
% Source for Figure~\ref{fig:lattice}. TikZ only; regenerable by compiling main.tex.
\begin{tikzpicture}[
  gd/.style={draw=rulegray!60},
  pt/.style={circle,fill=ink,inner sep=0.9pt},
  pt2/.style={circle,draw=ink,inner sep=0.9pt},
  pth/.style={->,accent,very thick},
  lab/.style={font=\footnotesize\sffamily,text=slate},
]
% 4x3 grid: x=0..4, y=0..3
\foreach \i in {0,...,4}
  \draw[gd] (\i*0.72,0) -- (\i*0.72,3*0.72);
\foreach \j in {0,...,3}
  \draw[gd] (0,\j*0.72) -- (4*0.72,\j*0.72);
% path R R U R U R U = (0,0)->(4,3)
\draw[pth] (0,0) -- (0.72,0);
\draw[pth] (0.72,0) -- (1.44,0);
\draw[pth] (1.44,0) -- (1.44,0.72);
\draw[pth] (1.44,0.72) -- (2.16,0.72);
\draw[pth] (2.16,0.72) -- (2.16,1.44);
\draw[pth] (2.16,1.44) -- (2.88,1.44);
\draw[pth] (2.88,1.44) -- (2.88,2.16);
\foreach \i in {0,...,4}
  \foreach \j in {0,...,3}
    \node[pt] at (\i*0.72,\j*0.72) {};
\node[lab,anchor=north] at (0,-0.1) {$(0,0)$};
\node[lab,anchor=west] at (4*0.72+0.08,3*0.72) {$(4,3)$};
\node[lab] at (5.6,2.0) {word:};
\node[lab] at (5.6,1.6) {\texttt{RRURURU}};
\node[lab] at (5.6,1.0) {4 R's, 3 U's};
\node[lab] at (5.6,0.6) {$\binom{7}{4}$ paths};
\node[lab] at (1.4,-0.6) {one path = one arrangement of the letters};
\end{tikzpicture}

\caption{The lattice-path bijection. Notice the word \texttt{RRURURU}
drawn as the path it encodes --- the picture \emph{is} the bijection,
not an illustration of it.}
\label{fig:lattice}
\end{figure}

\section{Stars and bars}\label{sec:stars}
\index{stars and bars}\index{multiset!counting}

\begin{theorem}
The equation $x_1 + x_2 + \cdots + x_k = n$ has $\binom{n+k-1}{k-1}$
solutions in nonnegative integers ($\binom{n-1}{k-1}$ in positive).
\end{theorem}

\begin{example}\label{ex:stars}
Solve $x_1 + x_2 + x_3 + x_4 = 12$ in nonnegative integers.
\par\smallskip
\textbf{Strategy.} Draw it before counting: write $n$ stars in a row,
place $k-1$ bars among the $n + k - 1$ symbol positions --- the bars cut
the stars into $k$ groups, and group sizes \emph{are} the solution.
\par\smallskip
\begin{solution}
Figure~\ref{fig:stars} shows one arrangement:
$\star\star\star\,\vert\,\star\star\,\vert\,\star\star\star\star\,\vert\,\star\star\star$
for $(3,2,4,3)$. Placing 3 bars among $12 + 3 = 15$ positions gives
$\binom{15}{3} = 455$ solutions.
\end{solution}
\end{example}

\begin{figure}[ht]\centering
% Figure: stars and bars --- 12 stars, 3 bars, 4 groups.
% Source for Figure~\ref{fig:stars}. TikZ only; regenerable by compiling main.tex.
\begin{tikzpicture}[
  st/.style={font=\small},
  bar/.style={draw=accent,very thick},
  lab/.style={font=\footnotesize\sffamily,text=slate},
]
% pattern: ***|**|****|***  = (3,2,4,3) sum 12, k=4 -> 3 bars
\foreach \i in {0,...,14}
{
  \ifnum\i=3 \draw[bar] (\i*0.5,-0.3) -- (\i*0.5,0.3); \else
  \ifnum\i=6 \draw[bar] (\i*0.5,-0.3) -- (\i*0.5,0.3); \else
  \ifnum\i=11 \draw[bar] (\i*0.5,-0.3) -- (\i*0.5,0.3); \else
    \node[st] at (\i*0.5,0) {$\star$};
  \fi\fi\fi
}
\node[lab] at (1.0,-0.9) {$x_1 = 3$};
\node[lab] at (3.0,-0.9) {$x_2 = 2$};
\node[lab] at (5.0,-0.9) {$x_3 = 4$};
\node[lab] at (6.9,-0.9) {$x_4 = 3$};
\node[lab] at (3.5,-1.6) {15 slots, choose 3 bar positions: $\binom{15}{3} = \binom{n+k-1}{k-1}$};
\end{tikzpicture}

\caption{Stars and bars for $x_1 + x_2 + x_3 + x_4 = 12$. Notice what a
``solution'' looks like physically: $k-1$ bars carving $n$ stars into
ordered piles --- order is free because the piles have names.}
\label{fig:stars}
\end{figure}

Positive solutions are the same picture with one star pre-committed per
group: put $1$ in each $x_i$ first, distribute $n - k$ stars freely.

\begin{exercise}[medium]\label{exr:stars}
Count solutions of $x_1 + x_2 + x_3 = 10$ with all $x_i \ge 2$. (Shift
variables first --- then stars and bars.)
\end{exercise}

\section{The pigeonhole principle}\label{sec:pigeon}
\index{pigeonhole principle}

\begin{theorem}[Pigeonhole]
If $n + 1$ objects go into $n$ boxes, some box gets at least two. More
generally, $N$ objects in $k$ boxes put at least $\lceil N/k \rceil$ in
one box.
\end{theorem}

Trivially true --- and load-bearing in olympiad proofs because the
\emph{boxes are a design choice}: partition the objects cleverly and
``two in one box'' becomes ``two objects sharing a property''.

\begin{example}\label{ex:pigeon}
Show that among any $51$ numbers from $\{1, \dots, 100\}$, two differ by
a factor of $2$.
\par\smallskip
\textbf{Strategy.} The boxes must make ``differing by a power of two''
mean ``same box''. Every integer is (odd part) $\times$ (a power of two);
numbers with the \emph{same odd part} differ only by powers of two.
\par\smallskip
\begin{solution}
Write each number as $2^e \cdot o$ with $o$ odd; there are exactly $50$
odd parts in $\{1, \dots, 100\}$, hence $50$ boxes
(Figure~\ref{fig:pigeon}). $51$ numbers force two into one class:
$a = 2^e o$ and $b = 2^f o$ with $e \ne f$, so one divides the other by
a factor of $2^{|e-f|} \ge 2$. Sharp bound: the $50$ odd numbers are a
$50$-element set with no such pair.
\end{solution}
\end{example}

\begin{figure}[ht]\centering
% Figure: pigeonhole for the 51-numbers example --- odd-core classes.
% Source for Figure~\ref{fig:pigeon}. TikZ only; regenerable by compiling main.tex.
\begin{tikzpicture}[
  box/.style={draw=rulegray,rounded corners=2pt,minimum height=8mm,inner sep=3pt},
  lab/.style={font=\footnotesize\sffamily,text=slate},
]
% 5 of the 50 classes: odd part 1,3,5,7,9 ... 99
\node[box,align=center,text width=1.55cm] at (0,0) {{$1,2,4,8,$\\$16,32,64$}};
\node[box,align=center,text width=1.55cm] at (2.1,0) {{$3,6,12,$\\$24,48,96$}};
\node[box,align=center,text width=1.55cm] at (4.2,0) {{$5,10,20,$\\$40,80$}};
\node[box,align=center,text width=1.55cm] at (6.3,0) {{$7,14,28,$\\$56$}};
\node[lab] at (7.9,0) {\dots};
\node[box,align=center,text width=1.55cm] at (9.4,0) {$99$};
\node[lab] at (0,0.9) {odd part $=1$};
\node[lab] at (2.1,0.9) {odd part $=3$};
\node[lab] at (4.2,0.9) {odd part $=5$};
\node[lab] at (6.3,0.9) {odd part $=7$};
\node[lab] at (9.4,0.9) {odd part $=99$};
\node[lab] at (4.7,-0.9) {50 classes cover $1..100$; two picks in one class differ by a factor of 2};
\node[lab] at (4.7,-1.4) {the all-odd 50-set shows 50 is the true threshold};
\end{tikzpicture}

\caption{The odd-part partition of $1..100$ into $50$ classes. Notice
the proof is in the box design: inside one class, two elements differ by
a pure power of $2$ --- the ``same box'' condition was engineered to be
exactly the conclusion.}
\label{fig:pigeon}
\end{figure}

\begin{checkpoint}[Predict before you run]
Listing~\ref{lst:count} checks the chapter's three closed forms against
brute force. Predict which count is \emph{not} a power of two before
reading the output.
\end{checkpoint}

\begin{lstlisting}[language=Python,caption={\texttt{count\_check.py}:
enumeration vs closed forms (full script in
\texttt{code/count\_check.py}).},label={lst:count}]
enum = sum(1 for w in combinations(range(m + k), m))
\end{lstlisting}
Real output:
\begin{lstlisting}[numbers=none]
subsets: n=4
  enum=16  closed=2^4=16
lattice paths: 3x2 grid
  enum=10  closed=C(5,3)=10
stars and bars: x1+x2+x3+x4 = 6, x >= 0
  enum=84  closed=C(9,3)=84
pigeonhole demo: 51 numbers from 1..100, two differ by 2x
  all-odd set [1, 3, 5, 7, 9]...[95, 97, 99]: no pair doubles -> 50 evades
  adding any even e makes e/2 or 2e odd-divisor-equal
\end{lstlisting}

\section{Committee proofs}\label{sec:committee}
\index{combinatorial proof}

The committee argument generalises WE~\ref{ex:committee}: an identity
between two formulas is proved by exhibiting \emph{one} set both sides
count. The discipline is identical to Chapter~2's proof hygiene --- name
the object, then let each side count it under a different strategy.
$\binom{n}{k} = \binom{n}{n-k}$: choosing a committee chooses its
non-members. Vandermonde $\sum_i \binom{m}{i}\binom{n}{k-i} =
\binom{m+n}{k}$: a $k$-committee from $m$ women + $n$ men splits by the
number of women chosen. Every algebraic identity in this chapter has a
committee behind it --- finding it is the exercise.

\begin{figure}[ht]\centering
% Figure: double counting — incidence table rows vs columns.
% Source for Figure~\ref{fig:doublecount}. TikZ only; regenerable by compiling main.tex.
\begin{tikzpicture}[
  cell/.style={draw,minimum size=7mm},
  tick/.style={font=\footnotesize,text=accent},
  hd/.style={font=\footnotesize\sffamily,text=slate},
  lab/.style={font=\footnotesize\sffamily,text=slate,align=center}
]
% 4 people x 3 committees incidence matrix
\foreach \c in {0,1,2} \node[hd] at (\c*0.75,0.45) {$C_{\c}$};
\foreach \r/\nm in {0/Ana,1/Bo,2/Cid,3/Dan} {
  \node[hd] at (-0.9,-\r*0.75) {\nm};
}
\foreach \r in {0,1,2,3} {\foreach \c in {0,1,2} {
  \node[cell] at (\c*0.75,-\r*0.75) {};
}}
% memberships: Ana in C0,C1; Bo in C0; Cid in C1,C2; Dan in C0,C2
\node[tick] at (0*0.75,-0*0.75) {\checkmark};
\node[tick] at (1*0.75,-0*0.75) {\checkmark};
\node[tick] at (0*0.75,-1*0.75) {\checkmark};
\node[tick] at (1*0.75,-2*0.75) {\checkmark};
\node[tick] at (2*0.75,-2*0.75) {\checkmark};
\node[tick] at (0*0.75,-3*0.75) {\checkmark};
\node[tick] at (2*0.75,-3*0.75) {\checkmark};
\node[lab] at (0.75,-3.5) {Rows sum $2{+}1{+}2{+}2 = 7$ checkmarks;\\
columns sum $3{+}2{+}2 = 7$};
\node[lab,align=center]{same set,\\
two orders\\
of counting};
\draw[->,thick,accent] (2.6,-1.0) -- (1.9,-1.0);
\end{tikzpicture}

\caption{Double counting in its rawest form: memberships as a
$\{\text{person}\}\times\{\text{committee}\}$ table. Row-sums and
column-sums both total $7$ --- and that tautology,
$\sum_i |C_i| = \sum_j |P_j|$, is the engine of every committee proof.}
\label{fig:doublecount}
\end{figure}

\begin{exercise}[medium]\label{exr:vandermonde}
Prove Vandermonde's identity by naming the set both sides count.
\end{exercise}

\section{Chapter summary}
\begin{table}[ht]\centering
\begin{tabularx}{\linewidth}{@{}lYY@{}}
\toprule
\tablehead\thead{Object} & \thead{Count} & \thead{Watch}\\
\midrule
ordered sequence, no repeats & $n!$, or $n!/(n-k)!$ for length $k$ & ties need their own rule\\
$k$-subset of $[n]$ & $\binom{n}{k}$ & unordered: no $k!$ factor\\
paths with $m$ R, $n$ U & $\binom{m+n}{m}$ & the word is the object\\
$x_1+\cdots+x_k = n$, $x_i \ge 0$ & $\binom{n+k-1}{k-1}$ & shift for lower bounds\\
same set, two routes & identity for free & name the set first\\
$N$ objects in $k$ boxes & one gets $\ge \lceil N/k \rceil$ & design the boxes\\
\bottomrule
\end{tabularx}
\caption{Counting in one page: pick the model (ordered? distinct?
bounded?), then let a bijection or a committee do the proving.}
\label{tab:count}
\end{table}

\section{Common mistakes}
\index{common mistakes}
\begin{itemize}
\item \textbf{Multiply where cases add.} ``Pick 1 of 3 or 1 of 4'' is
$7$ outcomes; $3 \cdot 4$ counts \emph{pairs}, a different object. Ask:
is the second choice conditional on the first?
\item \textbf{Ordered where unordered was meant.} Dealing a 5-card hand
by $52 \cdot 51 \cdots 48$ counts ordered hands; each real hand appears
$5!$ times. The fix ``divide by $k!$'' is honest here --- but only when
every target object was overcounted \emph{equally}.
\item \textbf{Ignoring ``indistinguishable''.} Stars and bars counts
named piles (the $x_i$ have labels). If the problem's piles are
unnamed, the same diagram overcounts --- partitions of $n$ have no
closed form, and pretending they do is the mistake.
\end{itemize}

\section{Additional exercises}
\begin{exercise}[easy]
Count 4-letter ``words'' from $\{a,b,c,d,e\}$: (a) any letters;
(b) distinct letters; (c) distinct with no two vowels adjacent.
\end{exercise}
\begin{exercise}[easy]
Evaluate $\binom{10}{3}$ and $\binom{10}{7}$ two ways each; which
computation is shorter and why does symmetry predict it?
\end{exercise}
\begin{exercise}[easy]
Count committees of 5 from 12 people if two specific people refuse to
serve together. (Two clean cases --- don't subtract blindly.)
\end{exercise}
\begin{exercise}[easy]
From a standard deck: how many 5-card hands are all hearts? All one
suit? Exactly 3 aces?
\end{exercise}
\begin{exercise}[medium]
Prove $\sum_{k=0}^{n} \binom{n}{k}^2 = \binom{2n}{n}$ by counting
$n$-subsets of $[2n]$ a second way.
\end{exercise}
\begin{exercise}[medium]
Count lattice paths $(0,0) \to (6,4)$ that pass through $(3,2)$ ---
then count those that avoid it. Which was easier to count, and why?
\end{exercise}
\begin{exercise}[medium]
Count nonnegative solutions of $x_1 + x_2 + x_3 + x_4 = 20$ with
$x_1 \le 6$. (Inclusion-exclusion on the bad side.)
\end{exercise}
\begin{exercise}[medium]
Prove: among any $6$ integers, two have difference divisible by $5$.
Generalise: $n+1$ integers give a pair with difference divisible by
$n$. What are the pigeons, what are the boxes?
\end{exercise}
\begin{exercise}[medium]
Show that any 10 distinct two-digit numbers contain two disjoint
subsets with equal sum. (Count subsets vs possible sums.)
\end{exercise}
\begin{exercise}[hard]
Prove $\binom{n}{0} + \binom{n}{2} + \binom{n}{4} + \cdots =
2^{n-1}$ by a committee argument \emph{and} by the binomial theorem
(evaluate at $(1,1)$ and $(1,-1)$).
\end{exercise}
\begin{exercise}[hard]
Derive the positive-solutions formula $\binom{n-1}{k-1}$ for
$x_1 + \cdots + x_k = n$ directly from the diagram (bars between stars,
at most one per gap), then reconcile it with the $x_i \ge 1$ shift.
\end{exercise}
\begin{exercise}[hard]
Extend \texttt{count\_check.py} to enumerate \emph{restricted}
stars-and-bars solutions ($x_i \le c_i$) for $k = 3$ and verify the
inclusion--exclusion count (the tool of Chapter~9) on 200 random cases.
\end{exercise}
\begin{exercise}[easy]
Count the arrangements of \textsc{banana} and of \textsc{bookkeeper} ---
three double letters in a row in the second one is not a special case,
just the same division-by-shuffles.
\end{exercise}
\begin{exercise}[medium]
A class of $12$ students splits into three groups of $4$ for projects.
Count the divisions --- and then recount if the groups are labelled
\emph{Red, Green, Blue}. (What does labelling multiply by?)
\end{exercise}
\begin{exercise}[hard]
Prove $\sum_{k=0}^{n} k\binom{n}{k} = n\,2^{n-1}$ two ways: once by
counting (a committee plus a named chair) and once by differentiating
$(1+x)^n$. Say which proof generalises to $\sum_k k^2\binom{n}{k}$ and why.
\end{exercise}

\section*{Further reading}
\begin{itemize}
\item R.~Hammack, \emph{Book of Proof}, 3rd ed., Ch.~3 ---
\url{https://richardhammack.github.io/BookOfProof/} --- free PDF;
counting, factorials and the binomial theorem.
\item E.~Lehman, F.~T.~Leighton, A.~R.~Meyer, \emph{Mathematics for
Computer Science}, Ch.~14--17 ---
\url{https://courses.csail.mit.edu/6.042/spring18/mcs.pdf} --- free PDF;
the sum/product rules, pigeonhole and counting-in-two-ways framing used
throughout this chapter.
\item Y.~Zhao, \emph{Bijections} ---
\url{https://yufeizhao.com/olympiad/bijections.pdf} --- the bijective
proof drill behind \S\ref{sec:biject}--\S\ref{sec:committee}.
\item Berkeley CS70 notes ---
\url{https://www.eecs70.org/} --- counting and discrete probability with
worked problems.
\item MIT OpenCourseWare, \emph{6.042J} (Fall 2010) ---
\url{https://ocw.mit.edu/6-042JF10} --- the counting lectures behind
\S\ref{sec:rules}--\S\ref{sec:pigeon}.
\item Cut-the-Knot, combinatorics ---
\url{https://www.cut-the-knot.org/} --- committee proofs of the binomial
identities in \S\ref{sec:binom}--\S\ref{sec:committee}.
\end{itemize}

\begin{rulebox}[Watch instead]
\begin{itemize}
\item \textbf{Spanning Tree --- ``What Is the Pigeonhole Principle?''},
8:23, \url{https://www.youtube.com/watch?v=B2A2pGrDG8I} --- the
principle plus three real applications (chessboard, planets, compression):
\S\ref{sec:pigeon}. (Verified 2026-10-11.)
\item \textbf{TrevTutor --- ``Discrete Math'' playlist}
\url{https://www.youtube.com/playlist?list=PLDDGPdw7e6Ag1EIznZ-m-qXu4XX3A0cIz}
--- its counting unit covers \S\ref{sec:rules}--\S\ref{sec:stars}
(permutations, combinations, stars and bars). (Verified 2026-10-11.)
\end{itemize}
\end{rulebox}

\chapternotes
The two-rules-first presentation and committee-proof discipline follow
Hammack's counting chapter~\cite{hammack}; 6.042J's counting-via-
bijection sections supplied the lattice-path and ``counting in two
ways'' framing~\cite{mcs}; CS70's counting notes carry the same
material with more identities~\cite{cs70}; Zhao's bijection and
counting-in-two-ways handouts are the olympiad continuation~\cite{zhao},
and Engel's counting chapter the hard-end drill~\cite{engel}. Contest
items powered by this chapter --- paths and stars-and-bars (AIME~2025~I~P3,
AIME~2024~I~P6, AIME~2024~II~P3), committee identities (AIME~2023~II~P10,
AIME~2020~I~P7), pigeonhole constructions (AIME~2025~II~P10--11,
AIME~2022~I~P6, AIME~2018~I~P9--10, AIME~2010~I~P10) and the
distribution-flavoured items AIME~2025~I~P10, AIME~2023~I~P11,
AIME~2021~I~P4, AIME~2020~I~P5, AMC~12~2024~B~P16, AMC~12~2023~A~P13 ---
are on the AoPS archive~\cite{aops-amc}.
